\documentclass[a4paper,11pt]{article}
\usepackage[]{graphicx}
\usepackage[]{xcolor}
\input{../../../Preambulos/preamble_esp}

%Comandos para personalizar la plantilla:
\renewcommand{\materia}{Optimización Matemática}
\renewcommand{\practica}{Ramificación y acotación}
\renewcommand{\curso}{2º grado en Inteligencia Artificial}



\decimalpoint

\usepackage{multicol}
\usepackage{multirow}

\usepackage{array}
\newcolumntype{C}[1]{>{\centering\let\newline\\\arraybackslash\hspace{0pt}}m{#1}}
\usepackage{colortbl}
\usepackage{etoolbox}

\DeclareMathOperator{\minim}{minimizar}
\DeclareMathOperator{\maxim}{maximizar}
\DeclareMathOperator{\subjectto}{sujeto\ a}
\DeclareMathOperator{\argmin}{argmin}
\DeclareMathOperator{\conv}{conv}
\DeclareMathOperator{\rango}{rango}

\renewcommand{\theenumi}{\upshape (\roman{enumi})}
\renewcommand{\labelenumi}{\upshape (\roman{enumi})} %Enumeraciones con n?meros romanos (nunca estar?n en cursiva) 

\newcommand{\sigF}{\mathcal F}
\newcommand{\sigG}{\mathcal G}

\newcommand{\R}{\mathbb R}
\newcommand{\N}{\mathbb N}
\newcommand{\Z}{\mathbb Z}

\newcommand{\x}{\boldsymbol{x}}
\newcommand{\xb}{\bar{\boldsymbol{x}}}
\newcommand{\y}{\boldsymbol{y}}
\newcommand{\z}{\boldsymbol{z}}
\renewcommand{\b}{\boldsymbol{b}}
\renewcommand{\c}{\boldsymbol{c}}
\renewcommand{\d}{\boldsymbol{d}}
\newcommand{\h}{\boldsymbol{h}}
\newcommand{\g}{\boldsymbol{g}}
\newcommand{\p}{\boldsymbol{p}}
\newcommand{\q}{\boldsymbol{q}}
\newcommand{\e}{\boldsymbol{e}}
\newcommand{\A}{\boldsymbol{A}}
\newcommand{\B}{\boldsymbol{B}}
\newcommand{\NN}{\boldsymbol{N}}
\newcommand{\Q}{\boldsymbol{Q}}
\newcommand{\I}{\boldsymbol{I}}
\renewcommand{\L}{\boldsymbol{L}}
\renewcommand{\H}{\boldsymbol{H}}
\newcommand{\V}{\boldsymbol{V}}
\renewcommand{\a}{\boldsymbol{a}}
\newcommand{\s}{\boldsymbol{s}}
\renewcommand{\v}{\boldsymbol{v}}
\renewcommand{\z}{\boldsymbol{z}}
\newcommand{\lambdab}{\boldsymbol{\lambda}}
\newcommand{\mub}{\boldsymbol{\mu}}
\newcommand{\zero}{\boldsymbol{0}}
\newcommand{\inter}[1]{\mathring{#1}}
\newcommand{\cl}[1]{\bar{#1}}

\newcommand{\ddd}[2]{\ensuremath{\frac{#1}{#2}}}

\newcommand{\fil}[2]{\ensuremath{{#2}^\text{\tiny f}_#1}}
\newcommand{\col}[2]{\ensuremath{{#2}^\text{\tiny c}_#1}}

\newcommand{\doble}[2]{\genfrac{}{}{0pt}{}{#1}{#2}}

\usepackage{relsize}
\renewcommand\RSsmallest{1pt}
%\renewcommand{\t}[1]{#1^{\text{{\tiny$\hspace{-0.08em}\mathsf{T}$}}}}
\renewcommand{\t}[1]{#1^{\text{{\relsize{-2}$\hspace{-0.01em}\mathsf{T}$}}}}



\newlength{\myheight}
\newsavebox{\mytext}
\savebox{\mytext}{%
	\Large\bfseries This is our title%
}
\settoheight{\myheight}{\usebox{\mytext}}
\newcommand{\rlogo}{$\vcenter{\hbox{\includegraphics[height=\myheight]{Rlogo}}}\,$}

%\newcommand{\solucion}[1]{#1}
%\newcommand{\solucion}[1]{}

\newcommand{\solucion}{1}

\usetikzlibrary{positioning,arrows.meta}

\newcommand{\subproblema}[8]{
	\node (L#1) at (#2,#3){};
	\node[draw,rectangle] (b#1) at (#2,#3) {
		\begin{minipage}[t]{2.4cm}
			\centering
			#7
		\end{minipage}
	};
	\node[draw, rectangle] (m#1)[above=0mm of b#1] {
		\begin{minipage}[t]{2.4cm}
			\centering
			#5 \\	#6
		\end{minipage}
	};
	\node[draw,rectangle,fill=black!30,minimum width=3.1cm] (t#1) [above=0mm of m#1] {
		\begin{tabular}{cr}
			$\boldsymbol{L_{#1}}$ & {\tiny $LB_{#1}=#4$}\\
		\end{tabular}
	};
	\node[draw,rectangle, rounded corners=2mm, minimum width=0.6cm] (l#1)[left=0.5mm of m#1] {#8};
}
\newcommand{\phan}{\phantom{LB}}



\IfFileExists{upquote.sty}{\usepackage{upquote}}{}
\begin{document}

\begin{example}
	Resuelve el siguiente problema de programación entera:
	\[
	\begin{array}{rlcl}
		(P1) \qquad \minim & -5x_1-8x_2 & & \\[1ex]
		\subjectto & \phantom{-1}x_1+\phantom{1}x_2 &\leq &6 \\
		& \phantom{-}5x_1+9x_2&\leq &45 \\
		& \multicolumn{1}{r}{\x}&\geq& 0 \\
		& \multicolumn{1}{r}{\x}&\in& \Z^2.
	\end{array}
	\]
\end{example}
\begin{resol}\phantom{p}

\begin{center}%
\shorthandoff{<>} %To avoidinterference with Babel
\setlength{\unitlength}{1cm}
\begin{tikzpicture}[>={Stealth[scale=3,width=4,length=4]},thick,every node/.style={minimum width=2.65cm,font=\footnotesize},lin/.style={black,thick},cut/.style={black!65,thin},set/.style={fill=black!20}]

 %\subproblema{j}{x}{y}{LB_j}{\xb}{\bar z}{UB-LB}{orden}

	\subproblema{1}{0}{0}{-\infty}{$\xb=(2.25,3.75)$}{$\bar z=-41.25$}{$LB=-41.25$}{1}
	\subproblema{2}{-2.5}{-3.5}{-41.25}{$\xb=(2,3.89)$}{$\bar z=-41.11$}{\phan}{2}
	\subproblema{3}{2.5}{-3.5}{-41.25}{$\xb=(3,3)$}{$\bar z=-39$}{$UB=-39$ \\ $\x^{best}=(3,3)$ \\ $LB=-41.11$}{3}
	
	\draw[->] (b1) --node[left=-0.3]{$x_1\leq 2$} (t2);
	\draw[->] (b1) --node[right=-0.3]{$x_1\geq 3$} (t3);
	
	\subproblema{4}{-5}{-7}{-41.11}{$\xb=(2,3)$}{$\bar z=-34$}{\phan}{4}
	\subproblema{5}{0}{-7}{-41.11}{$\xb=(1.8,4)$}{$\bar z=-41$}{$LB=-41$}{5}
	
	\draw[->] (b2) --node[left=-0.3]{$x_2\leq 3$} (t4);
	\draw[->] (b2) --node[right=-0.3]{$x_2\geq 4$} (t5);
	
	\subproblema{6}{-2.5}{-10.5}{-41}{$\xb=(1,4.44)$}{$\bar z=-40.56$}{\phan}{6}
	\subproblema{7}{2.5}{-10.5}{-41}{No}{Factible}{$LB=-40.56$}{7}

	\draw[->] (b5) --node[left=-0.3]{$x_1\leq 1$} (t6);
	\draw[->] (b5) --node[right=-0.3]{$x_1\geq 2$} (t7);
	
	\subproblema{8}{-5}{-14}{-40.56}{$\xb=(1,4)$}{$\bar z=-37$}{\phan}{8}
	\subproblema{9}{0}{-14}{-40.56}{$\xb=(0,5)$}{$\bar z=-40$}{$UB=-40$ \\ $\x^{best}=(0,5)$}{9}

	\draw[->] (b6) --node[left=-0.3]{$x_2\leq 4$} (t8);
	\draw[->] (b6) --node[right=-0.3]{$x_2\geq 5$} (t9);	

	%%%%%%Start grid
	%\draw[step=.5cm,gray, thin, dashed] (0,0) grid (6,6);
	%\foreach \x in {0,1,...,6}
	%\draw (\x, 0) node[anchor=north, darkgray, fill=white] {{\tiny $\x$}};
	%\foreach \y in {0,1,...,6}
	%\draw (0, \y) node[anchor=east, darkgray, fill=white] {{\tiny $\y$}};
	%%%%%%End grid
	\end{tikzpicture}
	\shorthandon{<>} %To avoidinterference with Babel
\end{center}
\end{resol}

\end{document}
